% Self-contained: compile twice with pdflatex; no external figures or .bib needed.
% Evidence: dense_knowledge_representation_complete.ipynb (saved experiment outputs).
\documentclass[10pt,letterpaper]{article}
\usepackage[margin=0.65in]{geometry}
\usepackage[T1]{fontenc}
\usepackage{lmodern,microtype,multicol}
\usepackage{amsmath,booktabs,titlesec}
\titleformat{\section}{\normalsize\bfseries}{}{0pt}{}
\titlespacing*{\section}{0pt}{5pt}{3pt}
\usepackage[hidelinks]{hyperref}
\hypersetup{pdftitle={Dense Knowledge Representation: CLIP versus Qwen3-VL Embeddings},
  pdfauthor={Longho Bernard Che},pdfsubject={CAP6938 Topic 1: MMEB failure analysis}}
\setlength{\columnsep}{0.25in}
\setlength{\parindent}{0pt}
\setlength{\parskip}{3pt}
\setlength{\tabcolsep}{4pt}
\newcommand{\heading}[1]{\par\addvspace{5pt}\noindent{\normalsize\bfseries #1}\par\nobreak}
\newcommand{\case}[1]{\par\addvspace{4pt}\noindent\textbf{#1}\par\nobreak}
\begin{document}
\begin{center}
{\Large\bfseries Dense Knowledge Representation:\\[2pt]
CLIP versus Qwen3-VL Embeddings}\par
\vspace{3pt}
Longho Bernard Che\quad $\vert$\quad CAP6938, Topic 1
\end{center}
\vspace{-8pt}
\begin{multicols}{2}
\heading{1. Experimental setup}
I compare CLIP and Qwen on 24 queries each from four MMEB datasets~\cite{mmeb}:
CIFAR-100 (classification), CIRR (image retrieval), TextVQA (answer retrieval),
and RefCOCO (object-crop retrieval). Seed 6938 controls sampling. Each query has
its annotated positive and 31 sampled negatives after text/path deduplication,
in shuffled order. Models rank identical pools; these are not official benchmark
scores.

\texttt{openai/clip-vit-base-patch32}~\cite{clip} independently encodes images
and text; my mixed-input adaptation adds their unit vectors equally and
renormalizes. \texttt{Qwen/Qwen3-VL-Embedding-2B}~\cite{qwen} jointly processes
image/text tokens and normalizes the last non-padding hidden state. Both embed
queries and candidates separately, ranking by cosine similarity. Query ablations
remove one modality while holding candidate embeddings fixed. The A100 40GB run and full
configurations are documented in the notebook.

\heading{2. Results}
\begin{center}
\small
\begin{tabular}{@{}lrrrr@{}}
\toprule
& \multicolumn{2}{c}{Accuracy (\%)} & \multicolumn{2}{c}{MRR}\\
Dataset & CLIP & Qwen & CLIP & Qwen\\
\midrule
CIFAR-100 & 4.2 & 58.3 & .203 & .696\\
CIRR      & 66.7 & 87.5 & .785 & .938\\
TextVQA   & 16.7 & 100.0 & .292 & 1.000\\
RefCOCO   & 91.7 & 100.0 & .939 & 1.000\\
\bottomrule
\end{tabular}
\end{center}
\vspace{-5pt}
Full-query accuracy and mean reciprocal rank (MRR) use 24 queries per row;
chance accuracy is 3.125\%. Overall accuracy is 44.8\% for CLIP and 86.5\% for
Qwen, with 40 CLIP-only failures and 13 shared failures.

Removing the generic classification instruction raises CLIP accuracy from 4.2\%
to 54.2\%, reducing Qwen's advantage to 4.2 percentage points. On CIRR, text-only
queries outperform full queries for both models (CLIP: 70.8\% versus 66.7\%;
Qwen: 91.7\% versus 87.5\%). These pools can reward the edit text without requiring
the reference image.

\heading{3. Five failure cases}
The notebook preserves the inspected images. Below, $F/I/T$ denotes the gold
target's rank for full, image-only, and text-only queries; lower is better.

\case{1. CIFAR-100:65 --- shared near-tie.}
For the gold label \emph{girl}, CLIP predicts \emph{man} ($F/I/T=2/1/2$) and
Qwen predicts \emph{woman} ($3/2/4$). CLIP's gold-minus-best-negative margin is
only $-0.0005$. Removing the instruction restores rank 1, identifying the added
text vector as the cause of the reversal. Qwen also worsens with text.
The coarse $32\times32$ portrait and ambiguous labels limit stronger claims.

\case{2. CIFAR-100:69 --- fusion loses recognition.}
For pink flowers labeled \emph{tulip}, CLIP predicts the literal dataset label
\texttt{cra} ($18/1/31$); Qwen retrieves \emph{tulip} ($1/1/13$). CLIP's gold
margin drops from $+0.0286$ for image alone to $-0.0816$ after fusion. Its image
embedding already recognizes the flower: adding generic text distorts the
ranking. Qwen retains the correct match through joint conditioning, although
candidate wording remains a confounder.

\case{3. CIRR:193 --- incomplete attribute binding.}
The reference shows four dogs in a car; the edit requests one dog in a field.
CLIP retrieves several dogs on a deck ($7/5/15$), missing count and setting.
Qwen retrieves one dog on a sofa ($2/12/1$), satisfying only count. Qwen's
text-only success disappears when the reference image returns, suggesting
competing appearance and edit cues in its pooled representation. Joint encoding
does not reliably bind every requested attribute.

\case{4. TextVQA:33 --- selecting relevant text.}
Asked the time on a smartwatch, CLIP predicts notification text \texttt{say thnx}
($6/2/12$); Qwen retrieves the small, partly glare-obscured \texttt{13:57}
($1/16/4$). Qwen succeeds only with both modalities, consistent with
question-conditioned selection of visual text. CLIP's independent image encoder
cannot adapt its features to the question. Different resolutions and OCR-related
training also contribute plausible explanations.

\case{5. RefCOCO:595 --- object versus scene.}
For ``white chair top right,'' CLIP retrieves diners ($5/5/2$), whereas Qwen
retrieves the chair crop ($1/4/1$). CLIP's global scene similarity can outweigh
the requested object because its visual features are not word-conditioned.
However, Qwen also succeeds from text alone: crop appearance can identify the
target. This supports language-to-crop matching without demonstrating spatial
reasoning. The metric measures retrieval, not localization IoU.

\heading{4. Conclusion}
Qwen wins more full-query rankings; TextVQA supports useful joint conditioning,
while CLIP's fusion rule causes substantial failures despite useful image
embeddings. Small samples, 32-way pools, and training/preprocessing differences prevent
causal attribution to architecture. Fusion-weight and harder-negative tests
could clarify these effects.

\begingroup
\footnotesize
\begin{thebibliography}{3}
\setlength{\itemsep}{0pt}
\setlength{\parskip}{0pt}
\bibitem{mmeb} Z. Jiang et al. \emph{VLM2Vec: Training Vision-Language Models for
Massive Multimodal Embedding Tasks}. 2024.
\href{https://arxiv.org/abs/2410.05160}{arXiv:2410.05160}.
\bibitem{clip} A. Radford et al. \emph{Learning Transferable Visual Models From
Natural Language Supervision}. 2021.
\href{https://arxiv.org/abs/2103.00020}{arXiv:2103.00020}.
\bibitem{qwen} M. Li et al. \emph{Qwen3-VL-Embedding and Qwen3-VL-Reranker:
A Unified Framework for State-of-the-Art Multimodal Retrieval and Ranking}. 2026.
\href{https://arxiv.org/abs/2601.04720}{arXiv:2601.04720}.
\end{thebibliography}
\endgroup
\end{multicols}
\end{document}
